\documentclass[11pt]{article}
\usepackage[a4paper,margin=1in]{geometry}
\usepackage[T1]{fontenc}
\usepackage{hyperref}
\usepackage{enumitem}
\usepackage{setspace}
\setstretch{1.2}

\begin{document}

\title{\textbf{Summary of Changes --- Manuscript \# Access-2026-23347}}
\author{}
\date{June 24, 2026}
\maketitle

\section*{Instructions for Reviewers}

Due to the complexity of the IEEE Access template (custom fonts, multi-column
layout, \texttt{\textbackslash PARstart}), the \texttt{latexdiff} tool could not
produce a compilable highlighted PDF. This document provides a structured
summary of every substantive change between the original submission and the
revised manuscript. All changes address reviewer comments or incorporate new
experimental data.

Page and line numbers refer to the revised manuscript PDF
(\texttt{llm\_annotation\_paper\_access\_revised.pdf}).

\section*{Major Changes}

\begin{enumerate}[leftmargin=*,label=\textbf{\arabic*.}]

  \item \textbf{Abstract (p.~1).} Fully rewritten with a clearer
  problem--method--findings--implication structure. Approximately 30\% shorter
  than the original. The new Abstract also summarises the expanded ten-model
  cross-model validation and the quality-threshold experiment.

  \item \textbf{Introduction, new paragraphs (p.~1--2).}
  \begin{itemize}
    \item Formal definition of \emph{visual anchoring}: ``the degree to which
    a category label can be recovered from a single cropped frame without
    recourse to discourse context, temporal information, or speaker intent.''
    \item Explicit motivation for agreement filtering: explains the intuitive
    rationale (dual-model agreement as a reliability heuristic) before
    discussing its failure modes.
  \end{itemize}

  \item \textbf{Methods, new Pipeline Overview subsection (p.~2).}
  Three-stage summary of the experimental pipeline (Annotation $\to$
  Filtering $\to$ Fine-tuning), with a compact enumerated list covering all
  design choices.

  \item \textbf{Methods, Table~1--2 merged (p.~3).}
  The original separate tables (Dataset Overview and Per-Class Counts) are
  merged into a single table with per-class bounding-box counts as sub-rows.

  \item \textbf{Methods, cross-model section expanded (p.~5).}
  The cross-model robustness check now describes ten MLLMs (expanded from
  five), spanning three model families (Qwen, LLaVA, Gemma) and parameter
  scales from 7B to 35B.

  \item \textbf{Results, expanded cross-model tables (p.~7--8).}
  \texttt{tab:crossmodel\_overall} and \texttt{tab:crossmodel\_intent} now
  include all ten models with verified accuracy numbers. All values are
  sourced from actual JSONL/CSV result files or execution logs.

  \item \textbf{Results, new quality-threshold experiment (p.~7).}
  New subsection reporting a targeted experiment with Qwen3.5-27B
  pseudo-labels on TeacherBehavior. Key finding: the critical annotation
  accuracy threshold at which pseudo-label fine-tuning surpasses the
  zero-shot baseline lies between 41\% and 50\%. Includes new
  \texttt{tab:quality\_threshold} and \texttt{fig:quality\_threshold}.

  \item \textbf{Discussion, Finding~3 rewritten (p.~9--10).}
  Updated to incorporate the quality-threshold evidence and the finding that
  the selective-routing gain scales with annotator quality.

  \item \textbf{Discussion, Limitations updated (p.~10--11).}
  The ten-model cross-model validation is now described as providing
  annotator-side external validity. The single-benchmark constraint is
  explicitly acknowledged.

  \item \textbf{Conclusions rewritten (p.~11--12).}
  Now includes: (a) the quality-threshold finding, (b) the capability
  boundary on guide/answer categories, and (c) a practical annotator
  selection guideline.

\end{enumerate}

\section*{Minor Changes (Reviewer~2, Comment~2.7)}

\begin{itemize}
  \item Long paragraphs in Introduction and Discussion split into shorter
  units.
  \item Technical terms defined at first use.
  \item Hedged compound sentences replaced by direct statements.
  \item Tables reordered and captions shortened to avoid repetition.
\end{itemize}

\section*{New Experimental Data Added}

\begin{itemize}
  \item Cross-model validation expanded from 5 to 10 models (Qwen2-VL-7B,
  LLaVA-1.5-7B, Qwen2.5-VL-7B, Gemma-3-27B, Qwen3.5-27B, Qwen3.5-35B-A3B,
  Qwen3.6-27B, Qwen3.6-35B-A3B, Gemma4-26B-A4B, Gemma4-31B).
  \item Qwen3.5-27B quality-threshold fine-tuning experiment on TeacherBehavior.
  \item All new data verified against result files on the experiment server.
\end{itemize}

\section*{Data Provenance}

All numerical values in the revised manuscript can be traced to specific
JSON/CSV/JSONL files in the accompanying GitHub repository
(\url{https://github.com/zhanglizhuo/llm-annotation}). Qwen3.6-27B values are
reconstructed from execution logs (result directory was lost before backup).

\end{document}
